\section{Experimental Setup}

All results in this paper come from reproducible artifacts, not from a human gold set. The main corpus consists of 100,032 normalized NVD records, 28,785 normalized GHSA records, and 8,066 aligned pairs. The current baseline discrepancy view is the 2026-07-15 input-integrity refresh of \texttt{build\_field\_discrepancies.py}.

The claim boundary matrix below records how each artifact is allowed to be used in the manuscript. It is part of the experimental setup because the current package deliberately mixes deterministic baselines, blank annotation templates, and silver labels without treating any of them as human gold.

\begin{table}[!htbp]
\centering
\setlength{\tabcolsep}{3pt}
\scriptsize
\caption{Generated table from 04\_experim.}
\label{tab:04-experim-table-1}
\begin{tabularx}{\textwidth}{>{\raggedright\arraybackslash}p{0.18\textwidth}>{\raggedright\arraybackslash}X>{\raggedright\arraybackslash}X>{\raggedright\arraybackslash}X>{\raggedright\arraybackslash}X}
\toprule
RQ & Artifact & Label source & Allowed claim & Not claimed \\
\midrule
RQ1 & discrepancy stats; RQ1 coverage summary & deterministic field rules & baseline discrepancy distribution across aligned NVD-GHSA fields & gold discrepancy labels or source correctness \\
RQ2 & typing diagnostics; blank annotation templates; references/CWE full-impact audits; fresh-CVE typing-stability cohort & deterministic rules, non-human Codex candidates, strict same-model dual-pass consensus, and blank human fields & rule-trigger/readiness diagnostics, bounded complete-impact candidate findings, and current-snapshot baseline stability & human-gold typing performance, independent-human agreement, candidate-profile gain, or production-rule validation \\
RQ3 severity & severity silver-v2 evidence and predictions & evidence-aware LLM silver labels & silver-label evidence-availability and baseline diagnostics & human-gold adjudication performance or final source truth \\
RQ3 affected\_versions & affected\_versions silver-v2 diagnostics; two CVE-disjoint frozen holdouts & evidence-aware LLM silver labels; strict agreement between blind Codex reviewers & development diagnostics, holdout agreement/coverage, task-separated method diagnostics, and protocol failure analysis & human-gold performance, independent-human validation, or semantic version-range adjudication \\
\bottomrule
\end{tabularx}
\end{table}

The field-coverage table separates raw NVD nonempty coverage from coverage inside the 8,066 aligned NVD-GHSA pairs. This table is descriptive: it explains which fields have one-sided missingness before discrepancy typing, not whether either source is correct.

\begin{table}[!htbp]
\centering
\setlength{\tabcolsep}{3pt}
\scriptsize
\caption{Generated table from 04\_experim.}
\label{tab:04-experim-table-2}
\begin{tabularx}{\textwidth}{>{\raggedright\arraybackslash}p{0.15\textwidth}>{\raggedright\arraybackslash}X>{\raggedright\arraybackslash}X>{\raggedright\arraybackslash}X>{\raggedright\arraybackslash}X>{\raggedright\arraybackslash}X>{\raggedright\arraybackslash}X}
\toprule
Field & NVD nonempty & GHSA nonempty & Both nonempty & NVD only & GHSA only & Both empty \\
\midrule
severity & 95549 & 8066 & 8033 & 0 & 33 & 0 \\
published/date & 100032 & 8066 & 8066 & 0 & 0 & 0 \\
references & 97988 & 8066 & 8037 & 0 & 29 & 0 \\
affected\_versions & 68784 & 8024 & 5667 & 27 & 2357 & 15 \\
cwe\_ids & 90738 & 7949 & 7870 & 60 & 79 & 57 \\
\bottomrule
\end{tabularx}
\end{table}

For RQ2, the analysis uses \path{results/rq2_discrepancy_typing/rq2_typing_diagnostics.json} as a non-gold diagnostic artifact. It aggregates deterministic \texttt{(field, status, rule note)} triggers over the 40,330 field instances in the 8,066 aligned pairs. The repository also contains a 300-instance stratified annotation seed at \path{data/annotations/rq2/discrepancy_typing_seed.jsonl} and \texttt{.csv}, with 60 sampled field instances per target field. A 20\% second-pass consistency review template is prepared under \path{data/annotations/rq2/consistency_review/}, with 60 rows and 12 rows per target field. The readiness diagnostic at \path{results/rq2_discrepancy_typing/rq2_sample_coverage.json} checks that all 17 nonzero field/status candidate strata and the top 12 full-corpus deterministic rule triggers are represented in the primary seed. These files are templates with blank primary and reviewer labels; they are not used to report human-gold accuracy, macro-F1, precision, recall, agreement, or Cohen's kappa.

A separate current-snapshot typing-stability cohort excludes 717 CVEs previously exposed in RQ2/RQ3 development, audits, or holdouts. It samples 250 rows per field with 1,250 globally unique CVEs, projects raw aligned values into inverse-order blind worklists, and seals six executable prediction profiles before either accepted reviewer output. Two Codex CLI pass roles use the same model/config but disjoint ephemeral session sets; strict consensus requires exact determinate labels, no low confidence, and no human-review request. The passes comprise 28 and 67 sessions rather than two persistent annotators. This design estimates non-human baseline stability only: all six sealed prediction columns are identical after exposure exclusion, so candidate-profile gain is structurally unidentifiable on this cohort.

To test whether a post-profile comparison is available, we later acquire an isolated official 2026 NVD feed and a commit-pinned reviewed-GHSA archive after the profile seal, without replacing the main corpus. Eligibility is frozen before review. The strict event-time tier requires both normalized publication timestamps to be later than the seal; only this tier may support a temporal-generalization claim. If it is unavailable, a separately named snapshot-external development tier may use previously unseen \texttt{CVE-2026-*} records under an availability-adaptive size rule. The acquisition finds no strict-event-time CVE but 5,948 snapshot-external single-GHSA CVEs, fixing the latter tier at 50 rows per field. Selection uses a 70\% proportional and 30\% equal-audit allocation over current-status strata, a fixed hash rank, and global CVE uniqueness. It does not use candidate predictions. Six profile columns, inverse-order blind worklists, source hashes, the Codex executable hash, and two pass identifiers are sealed before review. This design can make naturally occurring profile differences identifiable, but it remains development-only because collection time after the seal is not event time after the seal.

After the main comparison is unsealed, a post-selection evidence-secondary contract targets exactly the three rows where current and CWE-taxonomy profiles differ. Before two new reviewer files exist, it freezes nine URLs already listed by NVD or GHSA, their response bodies and metadata, the prior source and evaluation hashes, two inverse-order blind worklists, and literal allowed CWE path strings. Neither worklist contains a profile prediction or prior label. One reviewer judges the taxonomy relation before the concrete mechanism, while the other uses the reverse reasoning order. Strict consensus requires agreement on set relation, discrepancy label, taxonomy support, and whether the mapped weakness is supported by the frozen mechanism evidence. Two earlier protocol attempts are retained but excluded: v1 violates the confidence/uncertainty contract, and v2 violates the literal path interface. V3 uses fresh sessions and validates path values before writing. Because row selection follows the main comparison and all reviewers use the same model family, this stage can explain selected cases but cannot update the sealed 250-row evaluation, estimate method gain, or authorize promotion.

We then broaden the evidence view to all 50 sealed \texttt{cwe\_ids} rows, hiding the three known profile differences among 47 controls. The field-complete packet contains 33 exact sets, nine literal strict subsets, four overlap-non-subsets, and four disjoint pairs. It freezes at most three ranked NVD/GHSA references per row: 134 of 135 URL snapshots succeed. Two reverse-order worklists omit profile predictions, A/B labels, strict status, and profile-difference membership. V1 is rejected because it incorrectly forces every literal strict subset to \texttt{incomplete}; V2 permits evidence-dependent subset judgments but reaches merge with a nonliteral evidence quote. V3 moves rationale length, URL membership, literal-quote, duplicate-citation, and allowed-path checks before output is written. Its E/F roles use 20 disjoint ephemeral sessions. This design removes explicit difference-row selection from the reviewer input, but the CWE field and evidence contract are still chosen after A/B and profile unsealing. The result is therefore a post-hoc field-control diagnostic, not a replacement for the sealed main comparison.

One final evidence-secondary stage covers all 16 non-CWE rows among the 19 original A/B non-strict rows; the other three are excluded because the all-50 CWE audit already gives them strict field-complete decisions. Before retrieval, the stage fixes field counts (12 affected\_versions, two references, and two severity), a deterministic selector over at most six original NVD/GHSA URLs, reverse-order blind G/H worklists, citation gates, and development thresholds of 0.75 evidence availability, 0.40 strict secondary resolution, and 0.95 staged coverage. G/H see neither A/B or CWE decisions nor any profile prediction. The staged candidate may retain the 231 sealed strict rows and add only strict decisions from the two post-selected evidence stages. It is expressly diagnostic: neither passing nor failing may update the sealed profile evaluation, create human gold, or authorize promotion.

After A/B unsealing, a separate tiebreak protocol selects only the 103 non-strict rows from the original blind worklist. Reviewer C receives neither baseline predictions nor A/B outputs, uses the original prompt, and runs in ephemeral sessions disjoint from A/B. Before C output, the tiebreak manifest fixes a qualified-vote rule: at least two determinate, non-low-confidence votes without human-review requests must choose the same label. It also fixes advancement thresholds of 0.70 selected-row resolution and 0.975 combined cohort coverage. Because C shares the same model family, prompt, snapshot, and post-unsealing row selection, its output remains a non-human candidate diagnostic rather than independent-human adjudication.

After that gate fails, a second contract is fixed before new URL retrieval for all 37 remaining non-human candidate abstentions. It selects at most six URLs already present in each original NVD/GHSA reference context, freezes their extracted text/status records, and gives the same evidence-enriched row to reviewers D and E in opposite deterministic orders. D/E receive no baseline, A/B/C output, vote group, or candidate label. For affected\_versions, cwe\_ids, and references, secondary resolution requires strict D/E consensus and at least one successful frozen citation from each reviewer; severity may use its structured CVSS values. The pre-fetch gate requires 0.75 evidence availability, 0.40 secondary strict resolution, and 0.982 combined candidate coverage. This remains a post-unsealing, same-model-family, non-human development audit.

For real-human validation, the repository projects all 1,250 frozen rows into a separate source-bound three-stage packet. Review rows omit baseline status and notes, sampling strata, both Codex decisions, and their consensus. Two different real people must label each row independently before an author records and signs the resolution. An author-only scheduler prioritizes the 103 non-human reviewer disagreements and 216 baseline-versus-consensus disagreements but contains no labels; complete evaluation still requires all 1,250 rows. The validator binds source hashes and rejects source drift, partial pending content, identical reviewer IDs, unsigned resolutions, and any attempt to mark the packet itself as human gold. At the current snapshot all 1,250 rows are pending and none is signed.

The later 250-row post-profile cohort has an analogous full-cohort packet rather than a three-case form. It omits all six sealed predictions, A/B and E/F outputs, consensus, strata, and scheduler signals. Its author-only no-label queue places the 19 original non-strict rows first, then the three profile-difference diagnostics, 44 baseline-versus-consensus mismatches, and 184 completion rows. The blank packet passes ordinary integrity validation but fails both \texttt{--require-signed} and \texttt{--require-complete}; all 250 rows remain pending. File validation cannot establish that an identifier belongs to a real person or that the reviews were independently conducted, so external identity verification and separate human-gold promotion remain mandatory.

To isolate what human labels can change in the profile comparison, a label-free paired-outcome contract reads only the original sealed source and prediction files. It rejects hash drift, identifies all current-versus-CWE prediction differences, and enumerates every five-label assignment on those rows. Each assignment contributes a paired win, loss, or tie according to which sealed prediction matches. The assignment counts are logical cases, not probabilities or a prior over human decisions. A second label-free contract groups all six profiles by their complete prediction vectors and evaluates every profile pair with the conditional exact two-sided McNemar test at alpha 0.05. For the representative cross-vector pair, it enumerates whether each possible label makes current alone correct, the candidate alone correct, or both wrong. It also reports assumption-bound design sensitivities: the theoretical minimum correctness-discordant count for any rejection, conditional exact power under specified directional win probabilities, and future cohort sizes under a stationary observed prediction-difference rate. These diagnostics bound the possible full-cohort accuracy difference and the current cohort's paired-test capacity; they cannot establish absolute accuracy, macro-F1, construct validity, future disagreement prevalence, or the error distribution over equal-prediction rows, which remain part of the full real-person workflow.

Because the 250-row allocation is stratified by current status rather than sampled as a simple random draw, we separately run a prediction-only census over all five fields of all 5,948 snapshot-external eligible CVEs. The census binds the acquisition and profile code, exactly replays every sealed prediction on the selected 250 rows, and only then extends the same six deterministic profiles to 29,740 field instances. It creates neither a replacement cohort nor a reviewer worklist and reads no label, review, consensus, or evidence-secondary output. Its purpose is to measure deterministic profile disagreement availability in this revealed universe and to test whether the sampled 3/250 difference count can be used as a population-rate assumption. It cannot measure correctness, select a winning profile, validate a power model, or restore event-time eligibility.

After that census is revealed, two field-specific audits cover its complete 34-row prediction-difference union without treating the rows as a representative sample. The CWE audit retains all 29 \texttt{cwe\_ids} differences, freezes 85 ranked source-reference records (83 successful), and presents inverse-order worklists with current/candidate values removed. Every row visibly has a disjoint taxonomy relation and an official path, so selection blinding is explicitly incomplete. Strict E/F consensus requires exact agreement on relation, discrepancy label, taxonomy compatibility, and concrete mapping, plus literal frozen evidence. The references audit instead constructs profile-independent partitions over all raw URLs in the five-row references union. It restores NVD/GHSA membership only after review and reports two definitions separately: persistent underlying resource and frozen HTTP resource. Direct side and CVE fields are hidden, but URL text can expose CVE identity and alias structure. The first references seal is archived because its merge calls a nonexistent writer and produces no consensus; revision 2 re-seals the archived attempt, code, 26 URL probes, and fresh reviewer outputs before merge. Both audits are outcome-selected, same-snapshot, same-model-family diagnostics and cannot update the sealed 250-row evaluation or authorize promotion.

Because failure-mode inspection exposed a possible field-contract change, we additionally run a post-hoc cross-protocol stability gate before implementing either mechanism as a comparator. For severity, the fresh missing-score/vector projection is replayed on the 60 older primary rows and 12 older review rows using only their original structured field context. For affected\_versions, the audit compares input availability rather than inventing historical raw values: older seed rows can contain two empty projected range values while retaining package identity on one side, whereas fresh rows expose raw affected records. The advancement gate rejects a candidate if severity direction reverses across protocol generations, affected inputs are not comparable, or no signed human calibration establishes the intended construct. This gate can reject an unstable candidate but cannot validate either protocol.

We then run two sealed development calibrations without changing the comparator. Calibration v1 deterministically selects 60 fresh-holdout rows across four severity boundary patterns, one-sided unbounded affected claims, and unchanged controls. The two pass roles receive identical raw values and one prompt in reverse order through disjoint ephemeral sessions; exact/strict thresholds and per-stratum expected-label thresholds are fixed before review. After v1 exposes four residual constructs, calibration v2 freezes a refined prompt and gate before selecting 42 rows disjoint from every v1 row. V2 separates same-version from cross-version CVSS vectors, equal normalized ranges with different package names, singleton-versus-interval cases, and prerelease boundaries. Open affected-version strata are evaluated for strict reviewer stability rather than forced to share one expected label. For v2's sole non-strict row, a secondary packet freezes structured responses from the official GHSA advisory, fixing commit, and three linked XWiki Jira issues; two new disjoint sessions receive the same evidence, and the secondary gate additionally requires each reviewer to cite at least two frozen official URLs. After that conditional review disagrees, a deterministic projection audit checks component membership, product-to-artifact dependency edges, release-version policy, current and legacy catalog membership, source continuity across artifact migration, and prerelease ordering. Its first pass abstains before historical dependency edges are available; a second post-unsealing pass permits set comparison only after five product dependencies and the legacy/current source transition are frozen. We subsequently fix a generic graph schema and raw-data selector before fetching cross-case evidence. The selector does not read reviewer files: from the already unsealed v2 source it takes all eight two-sided affected\_versions rows with one subject per source, differing full artifact identifiers, and equal non-empty interval signatures. Each row must bind every endpoint using official boundary POMs, root Composer manifests, Go module manifests, or an explicit project/registry alias; missing evidence forces abstention. A second selector keeps the schema unchanged, excludes the XWiki development case, and takes all five remaining one-subject cross-artifact rows with non-equal signatures. Registry catalogs enumerate released versions before applying the fixed equal/subset/overlap relation map. Before running, advancement requires at least 4/5 projected rows and at least 4/5 candidates matching both sealed AI reviewers; reviewer files are read only after graph candidates are computed. After this no-go is observed, we freeze a separate Codex expert-contract candidate that chooses snapshot-extensional published-release semantics without claiming human approval. A label-independent one-to-many selector finds the sole InLong product row with two GHSA Maven components. Each component must bind through its official catalog and same-version InLong parent POM at all three boundaries; affected sets are evaluated separately, their heterogeneity is retained, and their union is compared with the NVD product set only when every edge and affected release is bound. Reviewer labels are consulted only after this candidate is computed. Both calibrations and all conditional post-unsealing audits remain non-human development evidence, not a confirmatory holdout or human gold.

The final graph stress test leaves the reviewed calibration source and selects from the complete aligned input. For each of NuGet, PyPI, and crates.io, it takes the minimum SHA-256-ranked CVE with one NVD product, exactly two GHSA packages, differing canonical component signatures, no more than three spans per claim, and parseable nonzero boundaries. This produces Oracle Database Server, LangChain, and Deno rows without consulting any reviewer or consensus artifact. Before registry retrieval, we fix edge classes for coordinated components, required dependencies, dependency constraints, parallel distributions, and unbound alternatives. A product-level union requires a deterministic total component-release-to-product-release map; an admissible dependency range is insufficient. Advancement requires at least 2/3 complete projections across at least two ecosystems. Fourteen official registry, version-metadata, dependency, and vendor responses are then frozen. This audit remains post-unsealing and non-human, but unlike the calibration-derived graph rows its upstream source is not conditioned on non-human consensus.

After this 0/3 result, we freeze a post-no-go recovery contract before retrieving any additional Deno evidence. The product domain is every stable official \path{denoland/deno} GitHub release from \texttt{1.41.3} through \texttt{2.3.2}, plus its immediate predecessor and successor. Every selected tag must contain one parseable exact \texttt{deno\_runtime} entry in its committed \texttt{Cargo.lock}; that runtime must exist in the frozen crates.io catalog, the complete sequence must be monotonic, and the two outer mappings must bound the claim \path{0.150.0 <= deno_runtime < 0.212.0}. Only then may the GHSA runtime interval be projected into product releases and unioned with the direct Deno intervals. An independent verifier reconstructs the dynamic source inventory and all sets from cached official responses. The fixed advancement threshold is 1/1, but passing remains a single-row, non-human development result.

We separately audit the three non-affected-version rows that remain unresolved after D/E. This cohort is explicitly selected from revealed outcomes, and protocol discovery inspects its evidence shape before v1; the run is therefore a targeted diagnostic with promotion disabled. The builder reconstructs exactly two \texttt{cwe\_ids} rows and one \texttt{references} row, removes prior labels and rationales from the worklist, and freezes 12 complete official responses. For \texttt{CVE-2024-8020}, a source gate parses the PyTorch Lightning \texttt{2.3.2} \texttt{post\_state} handler and requires the \path{POST /api/v1/state} route, request-body parsing, direct \texttt{body["state"]} indexing, and no local \texttt{try}. For \texttt{CVE-2023-4304}, the gate combines the official CWE-840 mapping restriction with the changed executable lines in the fixing patch and refuses to infer authorization semantics from name/email validation. For \texttt{CVE-2023-32187}, two predeclared reference-resource profiles retain exact successful HTTP resources or narrowly repair only the malformed suffix after a complete CVE Bugzilla identifier. The independent verifier reconstructs the cohort, evidence hashes, source AST, patch markers, and both reference-set relations without importing the analyzer.

We finally perform a post-hoc, read-only staged-frontier audit over the immutable A-E request logs and sealed merge summaries. Request identity is the ordered sample-ID payload; excess request attempts are counted per exact payload, while split retries may establish row completion without repairing attempt-level provenance. The audit separately counts request row-attempts, successful reviewer-row decisions, and token usage recorded on successful responses. It reports marginal candidate yield for A/B strict consensus, reviewer C, and evidence reviewers D/E against the unchanged 0.982 combined-coverage target. A predeclared stop rule rejects further same-model escalation only when both prior advancement gates fail, combined coverage remains below target, and the targeted residual audit promotes no row. This operational audit cannot create labels, impute missing error reasons or token usage, reduce the real-person review queue, or estimate accuracy.

After the D/E evidence stage leaves all 28 affected-version rows unresolved, we run a separate structural work-allocation audit. It reads the sealed blind rows and the previously frozen official-edge result, computes project-family and edge features, ranks repeated families, and only then attaches D/E labels as diagnostics. The fixed rules produce 16 families and select Mattermost ahead of LF Edge EVE and Hutool without using reviewer labels. The two fixed Mattermost rows share stable product/module spans but add Go pseudo-version and legacy-coordinate claims. Release-graph v1 stops before set analysis when the GitHub Releases API reaches its pagination cap; v2 also stops before manifests because two supplied boundary versions are not GitHub Release objects in the frozen prefix. An explicit post-failure v3 therefore tests a narrower 19-token official Git-tag manifest domain. Every tag must bind the current \path{server/v8} module; each pseudo-version must bind an exact commit, timestamp, module manifest, and determinate tag-to-commit ancestry for all 19 tokens; the legacy coordinate additionally requires a total exact-tag manifest mapping. Diverged histories and HTTP 404 records remain unresolved. Family advancement is fixed at 2/2, and a separate cache-only verifier reconstructs the full source inventory and every gate.

We next evaluate the second-ranked LF Edge EVE family under a separately frozen v1 mechanism contract. Protocol discovery before v1 is explicitly disclosed: it identifies a 207-token official Git-tag domain from \texttt{3.0.0} through \texttt{10.1.0}, the nested \path{pkg/pillar} Go module, the non-Go \path{pkg/vtpm} build component, and the paths changed by the two pseudo commits. This makes the run ineligible for candidate promotion or success-rate interpretation. To avoid hundreds of API comparisons, the analyzer freezes the official refs and a content-addressed Git pack containing every commit reachable from the 207 tags and pseudo commits plus the required tag, tree, manifest, and build objects. The cache-only verifier imports the pack into a new bare repository, validates object hashes, rebinds all tags to the official ref snapshot, and recomputes ancestry. Each row additionally requires an actual root Go module for the structured coordinate, advisory-to-repository and component binding, pseudo-path coherence, total main/LTS ancestry, and patched-anchor ancestry. Diverged histories, repository-name similarity, current UI prose, and timestamps cannot fill a missing edge. The family threshold remains 2/2 and all outputs are non-human diagnostics.

The third-ranked Hutool family is evaluated under a snapshot-extensional Maven v1 contract. Its disclosed pre-v1 discovery observes equal Maven Central catalogs for the aggregate \texttt{cn.hutool:hutool-all} artifact and the two advisory components, plus the aggregate structure at three critical versions; candidate promotion is therefore disabled regardless of the result. The stable domain is the exact intersection of the three cached catalogs after retaining only numeric \texttt{MAJOR.MINOR.PATCH} tokens and excluding five fixed milestone tokens. Every catalog must bind its requested coordinate and expose the same stable domain. Source \texttt{hutool-all} POMs at \texttt{5.8.19}, \texttt{5.8.21}, and \texttt{5.8.22} must declare same-parent-version dependencies on both components and a shade goal; the corresponding Maven Central aggregate JARs must contain compiled core and JSON classes. Catalog equality establishes total token correspondence over the frozen domain, while POM/JAR containment is claimed only at the three inspected anchors. Introduced-at-zero GHSA ranges are then evaluated over the cached component catalogs and unioned before comparison with the NVD product set. The family threshold remains 2/2, an independent verifier recomputes all sets from the hash-bound cache without network access, and a passing mechanism requires a later blind cohort rather than promotion on these revealed rows.

We subsequently reuse this frozen mechanism on a CVE-exposure-disjoint Hutool cohort from the same aligned snapshot. The exclusion parser projects only CVE identifiers from prior artifacts; baseline, reviewer, consensus, and candidate fields do not enter selection. The cohort builder scans all 8,066 matched rows, excludes the union of 1,967 previously exposed CVEs, and retains six Hutool rows that satisfy one of two predeclared routes: direct product-to-aggregate correspondence or product-to-aggregate-component correspondence. It seals the six rows before candidate computation. The mechanism predates this availability audit, but the cohort's availability and structural shape are observed in the same snapshot before sealing; we therefore preclude promotion and treat the run as a retrospective external application, not a future time holdout or human evaluation. The verifier independently reconstructs the exposure union, cohort, release sets, relations, and no-promotion gate from the authoritative aligned input and hash-bound v1 cache.

For the references candidate, we audit all 56 rows changed from representation discrepancy to incomplete. The builder recomputes every derived status, binds the exact changed-CVE set, and freezes 118 URL probes across 57 identity groups. All 56 rows enter both reviews; a 28-row automatic network-positive subset is not exempted. The reviewer worklist exposes neutral group IDs, original NVD/GHSA URLs, and frozen probes, but not the transformation, candidate identity, trigger stage, subset direction, automatic certificate, or expected label. Reviewer E inspects probes before URL structure, while reviewer F reverses that order. A protocol pilot is retained but excluded because its validator enforced a confidence rule absent from the sealed prompt. Revision 2 seals both complete role prompts and uses fresh ephemeral runs. Strict consensus requires matching determinate verdicts and group decisions with no additional-review request. The resulting narrower profile is selected after this audit and remains post-hoc, same-model-family, and non-human.

A separate blank human-review packet binds the same revision-2 seal, raw URLs, frozen probes, group order, and verdict-to-status contract for all 56 rows. It schedules a primary annotator, a different independent reviewer, and an author resolution/signoff. Because the non-human disagreement reflects two resource-identity constructs, each stage must explicitly select an underlying-content, frozen-HTTP, or written custom identity definition before recording group decisions. The 24 encoded-line rows are prioritized but not prelabelled. At the current snapshot all 56 rows are pending and none is signed; the packet and readiness artifact remain \texttt{label\_is\_human=false} and cannot support a human-gold claim.

For a non-human CWE candidate diagnostic, we additionally audit the complete 17-row impact set of \texttt{taxonomy\_v1}. The blind worklist and 34 current/candidate predictions are sealed before either reviewer file exists. Two fresh Codex runs receive only the blind CWE 4.20 context in isolated workspaces. Sixteen affected CVEs are disjoint from the 300-row primary seed. Strict consensus requires agreement on set relation, discrepancy label, and taxonomy-support verdict, with neither reviewer requesting additional review. This design tests the complete rule impact set but remains impact-selected, same-model-family, and ineligible for a human-gold claim.

After the first audit, its six unresolved rows and three strict candidate regressions form a post-hoc priority subset for an evidence-enhanced secondary audit. Before two additional reviewer files exist, a new builder freezes up to five ranked references listed by NVD or GHSA per row, the secondary prompt, the first-stage artifacts, and code hashes. The secondary blind worklist contains neither method prediction nor prior reviewer label. One new run reads frozen evidence before the CWE mapping; the other identifies the concrete mechanism first. A determinate decision must cite a literal substring from a successfully fetched snapshot. Secondary strict consensus requires agreement on relation, label, taxonomy support, and concrete-mechanism verdict. This stage diagnoses whether additional evidence resolves the selected boundary cases; because selection follows the first audit and evidence availability is uneven, it is not an independent confirmation.

For the severity adjudication prototype, the analysis uses the 80-sample severity seed and its evidence cache. The label distribution in the silver-v2 output is 75 factual\_conflict, 3 representation\_discrepancy, and 2 temporal\_discrepancy. The adjudicated-source distribution is 49 both, 26 NVD, 3 GHSA, and 2 abstain; confidence is 69 high and 11 medium.

For affected\_versions, the analysis uses the 100-sample manual-check set and the evidence-aware silver-v2 output. The silver label distribution is 47 factual\_conflict, 28 incomplete, 17 representation\_discrepancy, 4 equivalent, and 4 uncertain. The adjudicated-source distribution is 45 NVD, 35 both, 11 GHSA, 7 abstain, and 2 neither; confidence is 34 high, 64 medium, and 2 low.

We additionally froze a second 100-row affected\_versions holdout from the current 651 deterministic factual-conflict candidates. Selection used a fixed SHA-256 rank after excluding every CVE in the earlier 100-row development cohort, leaving 551 eligible CVEs. Evidence was collected in a dedicated cache; all 100 rows had at least one nonempty fetched-text record. A blind worklist exposed source values, context, and evidence but excluded previous annotations, predictions, and correctness fields. Predictions from 18 fixed methods were sealed before either reviewer decision file existed. Two Codex reviewers independently assigned both a discrepancy type and a supported source from the same blind worklist. The strict merge retained only determinate rows with exact joint agreement and valid URL-level support. These labels remain \texttt{label\_is\_human=false}: the design is development-disjoint and leakage-controlled, but it is neither human gold nor independent-human validation.

The preregistered v1 endpoint was source-selection accuracy over all strict joint-consensus rows. After unsealing, we found that this endpoint mixed factual-conflict source adjudication with non-conflict discrepancy typing. We therefore report the original endpoint unchanged and a clearly marked post-hoc split between factual-conflict and non-conflict rows.

We then froze a v2 cohort after excluding the 200 CVEs in both earlier affected\_versions cohorts, leaving 451 eligible deterministic conflict candidates before fixed-hash selection of 100 rows. Its dedicated snapshot contains 564 evidence records and usable fetched text for every row. Before either new reviewer file existed, we sealed 300 type predictions from three methods and 1,900 source predictions from 19 methods, together with method and protocol code hashes. V2 preregisters two endpoints: full discrepancy-type accuracy on strict type consensus, with method abstention counted as incorrect, and source accuracy only on rows with both strict factual-conflict type and strict source consensus. The reviewers cite literal frozen snippets through structured evidence claims. Every v2 decision and consensus row remains \texttt{label\_is\_human=false}.

After v2 was unsealed, we performed a clearly marked development-only leave-one-cohort-out diagnostic over the Phase D, v1, and v2 non-human labels. Gold-blind structured features exclude CVE identifiers, sample identifiers, discrepancy labels, and adjudicated sources. Balanced logistic regression and a shallow decision tree are trained on two cohorts and tested on the third. The advance gate requires a candidate to produce strictly more correct predictions than the best named comparator in every held-out cohort. Because the feature and model families were selected post-unseal, this diagnostic can reject a candidate but cannot confirm a method.

The severity diagnostic baselines are \texttt{prefer\_nvd}, \texttt{prefer\_ghsa}, \texttt{latest\_published}, and \texttt{evidence\_score\_baseline}. The affected\_versions silver-label diagnostic comparison uses the same fixed-source and latest-published baselines plus \path{version_token_support_baseline}, a simple fetched-text matcher that predicts support when evidence snippets mention at least one affected-version token from a side. The older URL-only LLM draft is not reused as an adjudication label set, and silver-v2 is evidence-aware but still not gold.

To prepare a human audit path for RQ3, the repository includes blank templates under \path{data/annotations/rq3/gold_audit/}: 80 severity rows and 100 affected\_versions development rows. These templates preserve the source evidence, source identifiers, evidence summaries, and silver-v2 annotations as provenance, but all \texttt{human\_audit} fields are still draft/blank. The guarded evaluator \path{experiments/rq3_adjudication/evaluate_rq3_human_audit.py} refuses the current templates and writes no gold-audit metrics until valid final human rows are present. Both CVE-disjoint holdout consensus artifacts also require real human signoff before they can support a human-gold claim.

The current paper artifact summary is written to \path{results/paper_cose/cose_artifact_tables.md} and \path{results/paper_cose/cose_artifact_tables.json}. Those files aggregate the reproducible counts used in the paper.
